\documentclass{beamer}
\usetheme{uic}
\usepackage{amsfonts,amsmath,oldgerm,algorithm,algpseudocode}
\usepackage[font=small,labelfont=bf]{caption} % Required for specifying captions to tables and figures

\newcommand{\hrefcol}[2]{\textcolor{uihteal}{\href{#1}{#2}}}
\newcommand{\testcolor}[1]{\colorbox{#1}{\textcolor{#1}{test}}~\texttt{#1}}
% Boxed placeholder for figures to be added later: {width}{height}{label}
\newcommand{\figplaceholder}[3]{\fbox{\parbox[c][#2][c]{\dimexpr#1-2\fboxsep-2\fboxrule\relax}{\centering\small\textcolor{steelgray}{[Figure: #3]}}}}

% Please see Section 18.1 of Beamer User Guide for all the options \usefonttheme provides
\usefonttheme[onlymath]{serif}
% \usefonttheme{serif} % use this if you would like Serif font throughout (and not just for math)


\titlebackground*{assets/uic_lockup_blue.pdf}

% NOTE 1: The asterisk splits the background image. This option is good
% for logo-based backgrounds. If you use an image based background
% it's recommended to not split it:

% \titlebackground{assets/uic_seo.jpg}

% NOTE 2: If you use a title background that does not have a logo, you might
% want to enable logo in the top left.
% To do that, simply comment out this line
\themecolor{lightnologo}

\title{AL-STARS: Agentic LLM-powered Synthetic Transportation Agent Response System}
\subtitle{Progress Update (ICT R27-293)}
% This can be adjusted accordingly for longer titles
\setlength{\titleboxwidth}{0.45\textwidth}

\author{Sybil Derrible, Isaac Salvador, Fatemeh Allahakbari\\University of Illinois Chicago}
\date{\today}

\begin{document}
\maketitle
\themecolor{light} % reverts to a logo based theme (if you disabled it for title page)

% enabled after title page creation (i.e. after \maketitle)
\footlinecolor{uicblue}


% ===================== Recap =====================

\begin{frame}{Agenda}
\begin{enumerate}
\item Project recap
\item Task 1: Literature review progress
\item Task 2: Architecture development kick-off
\item Next steps and discussion
\end{enumerate}
\end{frame}


\begin{frame}{Project Overview}
\begin{itemize}
\item \textbf{Title:} Agentic LLM-powered Synthetic Transportation Agent Response System (AL-STARS)
\item \textbf{Agency:} Illinois Department of Transportation (IDOT)
\item \textbf{PI:} Sybil Derrible, PhD, UIC
\item \textbf{Research Assistants:} Isaac Salvador, Fatemeh Allahakbari
\item \textbf{Duration:} May 2026 -- 2029 % TODO: confirm end date (workplan 8/15/2029 vs. kickoff 5/29/2029)
\item \textbf{Budget:} \$400,000 (\$300K federal, \$100K cost share)
\end{itemize}
\end{frame}


\begin{frame}{Motivation}
\framesubtitle{Household travel surveys (HTS) are getting harder to collect}
\begin{columns}
\begin{column}{0.55\textwidth}
\begin{itemize}
\item HTS underpin travel demand models and planning decisions
\item Sample rates: 1--3\% early on, $<$0.5\% by 2006, 0.04\% (CMAP 2019)
\item Low-income, transit-dependent, and minority groups undersampled
\item Self-selection and non-response bias
\end{itemize}
\end{column}
\begin{column}{0.45\textwidth}
\figplaceholder{\textwidth}{0.55\textheight}{HTS sample rate over time}
\end{column}
\end{columns}
\end{frame}


\begin{frame}{Why Agentic AI?}
\begin{itemize}
\item LLMs in transportation are mostly static pipelines (forecasting, imputation, classification)
\begin{itemize}
\item Tasks classical ML already handles
\end{itemize}
\item Agentic AI: LLMs as autonomous agents that reason, remember, and act
\item AL-STARS frames agents as \textbf{simulated travelers}
\begin{itemize}
\item Synthetic, one-to-one, realistic survey responses
\item Context-aware, explainable, behaviorally grounded
\end{itemize}
\item Goal: \emph{complement}, not replace, HTS collection
\end{itemize}
\end{frame}


\begin{frame}{Objectives and Tasks}
\begin{columns}
\begin{column}{0.5\textwidth}
\textbf{Deliverables}
\begin{itemize}
\item Agentic AI literature review
\item AL-STARS framework
\item AL-STARS web-based tool
\end{itemize}
\end{column}
\begin{column}{0.5\textwidth}
\textbf{Tasks}
\begin{enumerate}
\item Literature Review
\item Architecture Development
\item Web-based Application
\item Evaluation and Update
\item Dissemination
\end{enumerate}
\end{column}
\end{columns}
\end{frame}


\begin{frame}{Timeline}
\framesubtitle{We are here: October 2026}
\centering
\figplaceholder{0.9\textwidth}{0.6\textheight}{Gantt chart with ``we are here'' marker}
% TODO: reconcile Task 2/3 dates between workplan and kickoff deck
\end{frame}


% ===================== Task 1 =====================

\begin{frame}{Task 1: Literature Review}
\framesubtitle{Scope and method}
\begin{columns}
\begin{column}{0.55\textwidth}
\begin{itemize}
\item State-of-the-art and state-of-practice of agentic AI
\item Academic publications and industry applications
\item Search strategy: TODO
\item Papers screened / included: TODO
\end{itemize}
\end{column}
\begin{column}{0.45\textwidth}
\figplaceholder{\textwidth}{0.55\textheight}{Search/screening flow}
\end{column}
\end{columns}
\end{frame}


\begin{frame}{Organizing Framework}
\framesubtitle{Four-module agent taxonomy (Wang et al., 2024)}
\begin{columns}
\begin{column}{0.45\textwidth}
\begin{itemize}
\item \textbf{Profile:} identity and objectives
\item \textbf{Memory:} experience and context
\item \textbf{Planning:} goal decomposition
\item \textbf{Action:} interface with environment
\end{itemize}
\end{column}
\begin{column}{0.55\textwidth}
\figplaceholder{\textwidth}{0.6\textheight}{Agent modules diagram}
\end{column}
\end{columns}
\end{frame}


\begin{frame}{Agent Profile}
\framesubtitle{Persona generation}
\begin{columns}
\begin{column}{0.55\textwidth}
\begin{itemize}
\item Socio-demographic attributes via population synthesis
\item Psychological and social factors via role-playing
\item Key findings: TODO
\end{itemize}
\end{column}
\begin{column}{0.45\textwidth}
\figplaceholder{\textwidth}{0.55\textheight}{TODO}
\end{column}
\end{columns}
\end{frame}


\begin{frame}{Memory}
\begin{columns}
\begin{column}{0.55\textwidth}
\begin{itemize}
\item Short-term: prompt context
\item Long-term: external store via retrieval (RAG)
\item Behavioral consistency over time
\item Key findings: TODO
\end{itemize}
\end{column}
\begin{column}{0.45\textwidth}
\figplaceholder{\textwidth}{0.55\textheight}{TODO}
\end{column}
\end{columns}
\end{frame}


\begin{frame}{Planning and Reasoning}
\begin{columns}
\begin{column}{0.55\textwidth}
\begin{itemize}
\item Chain-of-Thought for simple reasoning
\item Multi-path reasoning (Tree of Thoughts, CoT-SC)
\item Feedback: environment signals, cross-agent interaction
\item Key findings: TODO
\end{itemize}
\end{column}
\begin{column}{0.45\textwidth}
\figplaceholder{\textwidth}{0.55\textheight}{TODO}
\end{column}
\end{columns}
\end{frame}


\begin{frame}{Action and Environment}
\begin{columns}
\begin{column}{0.55\textwidth}
\begin{itemize}
\item Text-based tools and APIs
\item Multi-agent settings
\item Key findings: TODO
\end{itemize}
\end{column}
\begin{column}{0.45\textwidth}
\figplaceholder{\textwidth}{0.55\textheight}{TODO}
\end{column}
\end{columns}
\end{frame}


\begin{frame}{LLMs as Survey Respondents}
\framesubtitle{Behavioral fidelity}
\begin{columns}
\begin{column}{0.55\textwidth}
\begin{itemize}
\item How well do LLM agents replicate human responses?
\item Evaluation metrics and benchmarks: TODO
\item Known limitations: TODO
\end{itemize}
\end{column}
\begin{column}{0.45\textwidth}
\figplaceholder{\textwidth}{0.55\textheight}{TODO}
\end{column}
\end{columns}
\end{frame}


\begin{frame}{Safety, Ethics, and Policy}
\begin{itemize}
\item Bias and representativeness: TODO
\item Privacy and data governance: TODO
\item Transparency and explainability: TODO
\item Deployment best practices: TODO
\end{itemize}
\end{frame}


\begin{frame}{Gaps and Remaining Work}
\begin{columns}
\begin{column}{0.5\textwidth}
\textbf{Gaps identified}
\begin{itemize}
\item TODO
\item TODO
\end{itemize}
\end{column}
\begin{column}{0.5\textwidth}
\textbf{Remaining (through Apr 2027)}
\begin{itemize}
\item TODO
\item Auxiliary review continues through project
\end{itemize}
\end{column}
\end{columns}
\end{frame}


% ===================== Task 2 =====================

\begin{frame}{Task 2: Architecture Development}
\framesubtitle{Objectives}
\begin{itemize}
\item LLM agents as simulated individuals from Illinois socio-demographic data
\item Four-module architecture: profile, memory, planning, action
\item Focus on behavioral fidelity over economic/heuristic assumptions
\item Success criteria: TODO
\end{itemize}
\end{frame}


\begin{frame}{Proposed AL-STARS Architecture}
\centering
\figplaceholder{0.9\textwidth}{0.65\textheight}{AL-STARS framework (Salvador et al., TRR 2026)}
\end{frame}


\begin{frame}{Synthetic Population Pipeline}
\begin{columns}
\begin{column}{0.5\textwidth}
\begin{itemize}
\item Census / ACS PUMS for socio-demographics
\item Travel data: CMAP My Daily Travel, NHTS
\item Population synthesis approach: TODO
\end{itemize}
\end{column}
\begin{column}{0.5\textwidth}
\figplaceholder{\textwidth}{0.6\textheight}{Data pipeline}
\end{column}
\end{columns}
\end{frame}


\begin{frame}{Survey Conversation Design}
\begin{columns}
\begin{column}{0.5\textwidth}
\begin{itemize}
\item Survey administrator agent $\leftrightarrow$ respondent agent
\item Data collection hidden from the respondent
\item Survey materials and question flow: TODO
\end{itemize}
\end{column}
\begin{column}{0.5\textwidth}
\figplaceholder{\textwidth}{0.6\textheight}{Example conversation}
\end{column}
\end{columns}
\end{frame}


\begin{frame}{Evaluation Plan}
\begin{itemize}
\item Statistical comparison against past HTS: TODO
\item Behavioral fidelity of individual responses: TODO
\item Scenario testing (e.g., extreme weather, pandemics)
\end{itemize}
\end{frame}


\begin{frame}{Infrastructure}
\begin{itemize}
\item LLM-agnostic framework
\item Model selection: TODO
\item GPU server for fine-tuning and inference
\end{itemize}
\end{frame}


% ===================== Wrap-up =====================

\begin{frame}{Upcoming Milestones}
\begin{itemize}
\item Nov 2026: Quarterly progress report
\item Early 2027: Task 3 (web application) begins % TODO: confirm Jan vs. Apr 2027
\item Apr 2027: Primary literature review complete
\item Jun 2027: Annual progress report and TRP meeting
\end{itemize}
\end{frame}


\begin{frame}{Discussion}
\begin{itemize}
\item Access to IDOT survey materials (e.g., CATI scripts)
\item Data access needs: TODO
\item Feedback on architecture direction
\end{itemize}
\end{frame}


\backmatter
\end{document}
